% !TEX TS-program = lualatex
% !TEX encoding = UTF-8

\documentclass[VesperaleOP.tex]{subfiles}

\ifcsname preamble@file\endcsname
  \setcounter{page}{\getpagerefnumber{M-vop09_tempus_paschale}}
\fi

\begin{document}

%%%% IN DIE SANCTO PASCHÆ
\festum{}{In die Sancto Paschæ}{I classis cum octava I classis}{In die Sancto Paschæ}{1}

\titulumrubrica{Ad Vesperas}
\cantus{P0F1VA1}{Super Ps.}
\rubrica{Finita antiphona statim dicitur :}
\cantus{P0VR}{Ant.}

\cantus{P0F1VAM}{Ad Magn.}
\oratio{Deus, qui hodiérna die per Unigénitum tuum æternitátis nobis áditum, devícta morte reserásti:\psstar{} vota nostra, quæ præveniéndo aspíras,\pscross{}étiam adiuvándo proséquere.\psstar{} Per eúndem Dóminum.}

\titulumrubrica{Ad Processionem}
\rubrica{Hoc triduo terminata ad Vesperas orationem sequitur :}
\cantus{P0F1VA2}{Ant.}

\versiculus{Dícite in natiónibus, allelúia.\pscross{}}{Quia Dóminus regnávit a ligno, allelúia.\pscross{}}
\oratio{Deus, qui pro nobis Fílium tuum Crucis patíbulum subíre voluísti, ut inimíci a nobis expélleres potestátem,\pscross{}concéde nobis fámulis tuis ut resurrectiónis grátiam consequámur.\psstar{} Per eúndem Christum.}

\cantus{P0F1VA3}{Ant.}

\versiculus{Ora pro nobis sancta Dei Génetrix, allelúia.\pscross{}}{Ut digni efficiámur promissiónibus Christi, allelúia.\pscross{}}
\oratio{Grátiam tuam, quǽsumus, Dómine, méntibus nostris infúnde:\psstar{} ut qui, Angelo nuntiánte, Christi Fílii tui incarnatiónem cognóvimus,\pscross{}per passiónem eius et crucem ad resurrectiónis glóriam perducámur.\psstar{} Per eúndem Christum.}

\dominusvobiscumversiculus
\cantus{ORB8}{}

\lineamagna

%%%% INFRA OCTAVAM PASCHÆ
\festum{}{Feria secunda}{}{Feria II infra octavam Paschæ}{2}
\rubrica{Omnia ut heri, præter Antiphonam ad Magnificat et Orationem :}
\cantus{P0F2VAM}{Ad Magn.}
\oratio{Deus, qui solemnitáte pascháli mundo remédia contulísti;\pscross{}pópulum tuum, quǽsumus, cælésti dono proséquere:\psstar{} ut et perféctam libertátem cónsequi mereátur,\pscross{}et ad vitam profíciat sempitérnam.\psstar{} Per.}

\festum{}{Feria tertia}{}{Feria III infra octavam Paschæ}{2}
\rubrica{Omnia ut in die sancto Paschæ, præter Antiphonam ad Magnificat et Orationem :}
\cantus{P0F3VAM}{Ad Magn.}
\oratio{Deus, qui Ecclésiam tuam novo semper fetu multíplicas:\psstar{} concéde fámulis tuis ut sacraméntum vivéndo téneant,\pscross{}quod fide percepérunt.\psstar{} Per.}

\festum{}{Feria quarta}{}{Feria IV infra octavam Paschæ}{3}
\cantus{P0F4VAM}{Ad Magn.}
\oratio{Deus, qui nos resurrectiónis Domínicæ ánnua solemnitáte lætíficas:\psstar{} concéde propítius ut per temporália festa, quæ ágimus,\pscross{}perveníre ad gáudia ætérna mereámur.\psstar{} Per eúndem.}

\rubrica{Hodie et deinceps usque ad feriam VI inclusive fit tantum memoria de Cruce.}
\cantus{P0VA}{Ad mem.}

\versiculus{Dícite in natiónibus, allelúia.\pscross{}}{Quia Dóminus regnávit a ligno, allelúia.\pscross{}}
\oratio{Deus, qui pro nobis Fílium tuum Crucis patíbulum subíre voluísti, ut inimíci a nobis expélleres potestátem,\pscross{}concéde nobis fámulis tuis ut resurrectiónis grátiam consequámur.\psstar{} Per eúndem.}

\festum{}{Feria quinta}{}{Feria V infra octavam Paschæ}{3}
\cantus{P0F5VAM}{Ad Magn.}
\oratio{Deus, qui diversitátem géntium in confessióne tui nóminis adunásti:\psstar{} da ut renátis fonte baptísmatis una sit fides méntium,\pscross{}et píetas actiónum.\psstar{} Per.}

\festum{}{Feria sexta}{}{Feria VI infra octavam Paschæ}{3}
\cantus{P0F6VAM}{Ad Magn.}
\oratio{Omnípotens sempitérne Deus, qui paschále sacraméntum in reconciliatiónis humánæ fœ́dere contulísti:\psstar{} da méntibus nostris ut, quod professióne celebrámus,\pscross{}imitémur efféctu.\psstar{} Per.}
\lineamagna
\end{document}